\section*{Capítulo \ref{ch:b2:cell-signalling} --- Mensajeros químicos y transducción de señales}

\begin{solution}{exo:b2:cell-signalling:1}
Insulina: endocrina, hidrosoluble (péptido). Cortisol: endocrino,
liposoluble. Acetilcolina en una sinapsis: sináptica, hidrosoluble.
Histamina en una herida: paracrina, hidrosoluble. Estradiol:
endocrino, liposoluble. Adrenalina de la glándula suprarrenal: endocrina,
hidrosoluble. Óxido nítrico: paracrino, liposoluble (un gas que
atraviesa las membranas).
\end{solution}

\begin{solution}{exo:b2:cell-signalling:2}
La adrenalina se une al receptor (apagado: disociación, internalización
del receptor); el receptor activa la \omterm{prop:b2:cell-signalling:gpcr}{proteína G}, GDP por GTP (apagado:
hidrólisis del GTP en segundos); la \omterm{prop:b2:cell-signalling:gpcr}{proteína G} activa la adenilato ciclasa,
que fabrica AMPc (apagado: fosfodiesterasa); el AMPc activa la quinasa A
(apagado: pérdida del AMPc); la quinasa A fosforila la fosforilasa quinasa,
que fosforila la glucógeno fosforilasa (apagado: fosfatasas); la
fosforilasa libera glucosa-1-fosfato del glucógeno, que se transforma y
se exporta como glucosa.
\end{solution}

\begin{solution}{exo:b2:cell-signalling:3}
Los receptores de superficie unen mensajeros hidrosolubles en la membrana y
actúan en segundos mediante \omterm{prop:b2:cell-signalling:gpcr}{segundos mensajeros} y quinasas, cambiando la
actividad de proteínas ya existentes. Los \omterm{prop:b2:cell-signalling:nuclear}{receptores nucleares} unen
mensajeros liposolubles dentro de la célula y actúan en horas como
\omterm{prop:b2:cell-differentiation:transcriptional}{factores de transcripción}, cambiando qué proteínas se fabrican.
\end{solution}

\begin{solution}{exo:b2:cell-signalling:4}
\omterm{prop:b2:cell-signalling:gpcr}{AMP cíclico}: lo fabrica la adenilato ciclasa y lo destruye la
fosfodiesterasa. Calcio: entra por canales o sale del retículo
endoplásmico en respuesta al inositol trifosfato, y lo retiran las
bombas. Inositol trifosfato (con el diacilglicerol): lo fabrica la
fosfolipasa C a partir de un lípido de membrana, y lo retiran las fosfatasas.
\end{solution}

\begin{solution}{exo:b2:cell-signalling:5}
$\theta = [L]/(K_d + [L])$: $0.09$, $0.5$, $0.91$, $0.99$. Duplicar los
receptores deja la fracción sin cambios y duplica el número de receptores
ocupados --- y, por tanto, la respuesta.
\end{solution}

\begin{solution}{exo:b2:cell-signalling:6}
$50\times 200\times 100\times 300 = 3\times 10^{8}$. Cien receptores
ocupados dan $3\times 10^{10}$ moléculas de producto por segundo --- un
límite superior, puesto que las etapas se saturan.
\end{solution}

\begin{solution}{exo:b2:cell-signalling:7}
La \omterm{prop:b2:cell-signalling:gpcr}{proteína G} bloqueada mantiene activa la ciclasa, el AMPc sigue alto, la
quinasa A mantiene fosforilado y abierto el canal de cloruro de las
células intestinales, se secreta cloruro, y el sodio y el agua lo siguen
hacia la luz intestinal --- litros al día. El bloqueo es una modificación
covalente de la \omterm{prop:b2:cell-signalling:gpcr}{proteína G}, que solo se deshace cuando las propias células
se renuevan, días después.
\end{solution}

\begin{solution}{exo:b2:cell-signalling:8}
$0.9\times 10^{-6}\,\text{mol/L}\times 2\times 10^{-12}\,\text{L} =
1.8\times 10^{-18}$ mol, $1.1\times 10^{6}$ iones (en realidad más, pues
los tampones fijan la mayor parte de lo que entra). Un solo canal a
$10^{6}$ por segundo tardaría un segundo; mil canales lo hacen en un
milisegundo, que es la escala de tiempo que necesita una señal.
\end{solution}

\begin{solution}{exo:b2:cell-signalling:9}
Las dos vías convergen en las mismas enzimas: la quinasa A del glucagón
fosforila la glucógeno sintasa (que se desactiva) y la fosforilasa quinasa
(que se activa); la cascada de la insulina activa la fosfatasa que retira
los fosfatos. El estado de las enzimas lo fija el equilibrio entre la
actividad de las quinasas y la de las fosfatasas, de modo que lo que
cuenta es el cociente entre las dos \omterm{def:b2:cell-signalling:kinds}{hormonas}; con ambas altas, suele ganar
la fosfatasa y el glucógeno se almacena.
\end{solution}

\begin{solution}{exo:b2:cell-signalling:10}
Adrenalina: el AMPc difunde a través de una célula de \qty{20}{\micro m}
en menos de un segundo ($x^{2}/2D$ con $D \approx \qty{3e-10}{m^2/s}$), y
todas las enzimas de la cascada existen ya --- segundos. Cortisol: un
transcrito tarda \qty{60}{s}, una proteína \qty{100}{s}; con unas pocas
decenas de ARNm, cada uno con una docena de ribosomas, la célula fabrica
unos cientos de moléculas de enzima por minuto, de modo que los miles
necesarios para un efecto medible tardan horas, tras los minutos que
necesita el cortisol para entrar y llegar al ADN.
\end{solution}

\begin{solution}{exo:b2:cell-signalling:11}
El receptor emite señales sin cesar: la cascada de las MAP quinasas empuja
a la célula a dividirse sin ningún factor de crecimiento, y la célula
ignora la ausencia de la señal que normalmente la frena. Como una sola
\omterm{def:b2:genome-diversification:mutation}{mutación} elimina un control de la división, estos receptores (y las
quinasas situadas aguas abajo) están entre los oncogenes más frecuentes.
Un fármaco que ocupa el sitio del ATP de la quinasa detiene las
fosforilaciones y silencia el receptor --- el principio de varios
fármacos contra el cáncer.
\end{solution}

\begin{solution}{exo:b2:cell-signalling:12}
Endocrino: llega a todas las células en un minuto, es barato por célula,
no puede dirigirse a una sola célula y actúa de minutos a días ---
imprescindible para los estados del cuerpo entero (ayuno, crecimiento,
reproducción). Nervioso: un milisegundo, una sola diana, un cableado caro,
actúa durante milisegundos --- imprescindible para el movimiento y los
reflejos rápidos. Ambos se usan para el mismo mensaje cuando se necesitan
a la vez velocidad y alcance: los nervios simpáticos y la médula
suprarrenal liberan noradrenalina o adrenalina durante un susto.
\end{solution}

\begin{solution}{pb:b2:cell-signalling:1}
\textbf{1.} $1\,\text{nmol}/5\,\text{L} = \qty{0.2}{nmol/L}$.
\textbf{2.} $\theta = 0.2/1.2 = 0.17$: unos 3300 receptores ocupados.
\textbf{3.} Tres semividas: \qty{0.025}{nmol/L}; $\theta = 0.024$.
\textbf{4.} De medio minuto a un minuto, el tiempo que tarda la \omterm{def:b2:blood-circulation:blood}{sangre} en
dar la vuelta; el nervio entrega su transmisor directamente al hígado en
menos de un segundo.
\textbf{5.} La misma fracción, pero $33\,000$ receptores ocupados: una
respuesta diez veces mayor.
\textbf{6.} A \qty{1}{nmol/L}, un receptor con $K_d = \qty{1}{\micro
mol/L}$ está ocupado en una milésima: no hay respuesta. Ajustar la $K_d$ al
intervalo circulante sitúa la parte empinada de la curva de ocupación
donde la \omterm{def:b2:cell-signalling:kinds}{hormona} varía realmente.
\textbf{7.} $20\times 100\times 1\times 50\times 50\times 1000 =
5\times 10^{9}$ moléculas de glucosa por segundo por receptor ocupado.
\textbf{8.} $3300\times 20\times 100 = 6.6\times 10^{6}$ AMPc en el
primer segundo; en \qty{5e-12}{L}, $1.1\times 10^{-17}$ mol, unos
\qty{2}{\micro mol/L}.
\textbf{9.} $3300\times 5\times 10^{9} = 1.7\times 10^{13}$ por segundo, $2\times
10^{13}$ por minuto por célula.
\textbf{10.} $2\times 10^{11}$ células $\times 10^{15} = 2\times
10^{26}$ moléculas por minuto, 330 mol, \qty{6e4}{g} por minuto ---
diez mil veces los \qty{5}{g/min} observados. Las ganancias solo se
multiplican mientras ninguna etapa está saturada; en realidad la
fosforilasa está limitada por su sustrato y por las fosfatasas, y la
exportación de glucosa por los transportadores. La ganancia de la cascada
se invierte en rapidez y sensibilidad, no en producción.
\textbf{11.} A \qty{5}{g/min}, veinte minutos; el susto termina en unos
minutos y los interruptores de apagado detienen la cascada mucho antes.
\textbf{12.} La hidrólisis del GTP por la \omterm{prop:b2:cell-signalling:gpcr}{proteína G} (segundos); la
fosfodiesterasa sobre el AMPc (segundos); las fosfatasas sobre las
enzimas fosforiladas (de segundos a minutos); la desensibilización de los
receptores (minutos). Los más rápidos son el reloj propio de la \omterm{prop:b2:cell-signalling:gpcr}{proteína
G} y la fosfodiesterasa.
\textbf{13.} El AMPc no se destruye: se acumula y persiste, la quinasa
sigue activa, y la producción de glucosa es mayor y dura más --- el
estado de alerta que da el café.
\textbf{14.} $4000/50 = \qty{80}{s}$.
\textbf{15.} $600/5 = \qty{120}{s}$ por enzima por ribosoma; 30 por
ribosoma por hora, 600 por ARNm por hora.
\textbf{16.} $100\times 600 = 60\,000$ enzimas por hora; $10^{5}$ al cabo
de unas \qty{1.7}{h}.
\textbf{17.} De diez a veinte minutos para que el cortisol entre, se una y
llegue a sus genes, una hora antes de que los primeros ARNm sean
numerosos y se traduzcan, y luego dos horas de acumulación: unas cuatro
horas.
\textbf{18.} Una cascada cambia la actividad de enzimas que ya existen; la
función del cortisol es cambiar qué enzimas existen, durante días, cosa que
solo puede hacer la transcripción. Un susto no puede esperar cuatro horas
a que haya enzimas nuevas.
\textbf{19.} Más receptores significa más receptores ocupados a
cualquier concentración de adrenalina: la fracción de ocupación no
cambia, pero la respuesta a cada concentración es mayor --- la curva se
amplía en vertical, no se desplaza. Una \omterm{def:b2:cell-signalling:kinds}{hormona} lenta fija la ganancia de
una rápida.
\textbf{20.} La glucosa entra en la célula $\beta$ y se metaboliza; el ATP
sube y cierra un canal de potasio sensible al ATP; la membrana se
despolariza; se abren canales de calcio dependientes de voltaje; el calcio
desencadena la exocitosis de los gránulos de insulina --- el metabolismo
acoplado a la secreción mediante una etapa eléctrica.
\textbf{21.} Los dos receptores inician dos vías que convergen en las
mismas enzimas: la quinasa A del glucagón las fosforila, la cascada de la
insulina activa la fosfatasa que las desfosforila; el estado de las
enzimas, y por tanto el sentido del metabolismo del glucógeno, sigue el
cociente entre las dos \omterm{def:b2:cell-signalling:kinds}{hormonas}.
\textbf{22.} Glucemia en ayunas alta (el hígado, sin freno, sigue
fabricando glucosa), glucemia tras la comida muy alta (sin captación por
el músculo y el tejido adiposo); sin insulina, las reservas de grasa
liberan ácidos grasos, que el hígado convierte en cetoácidos ---
cetoacidosis.
\textbf{23.} El músculo en actividad capta glucosa sin insulina, de modo
que en una carrera larga la glucosa baja mientras la insulina inyectada,
sin regulación, sigue frenando el hígado: hipoglucemia y colapso.
\textbf{24.} Bucle de la glucosa: como sensores, las células $\beta$ y $\alpha$;
dos \omterm{def:b2:cell-signalling:kinds}{hormonas} efectoras para los dos sentidos; una ganancia alta que
mantiene el nivel con un margen de una quinta parte; minutos para actuar.
\omterm{def:b2:blood-pressure:baroreflex}{Barorreflejo}: receptores de estiramiento; el \omterm{def:b2:heart:anatomy}{corazón} y los vasos como
efectores; ganancia de alrededor de cuatro; un segundo para actuar.
\textbf{25.} Ocupación de $0.17$ tras la liberación; ganancia de $5\times 10^{9}$
por segundo por receptor ocupado; producción del hígado de
\qty{5}{g/min} en realidad; el cortisol tarda unas cuatro horas.
\end{solution}
